\documentclass[10pt,a4paper]{amsart}
\usepackage[margin=19mm]{geometry}
\usepackage{amsmath,amssymb,mathtools}
\usepackage[scr=rsfs,cal=euler]{mathalfa}
\usepackage{xfrac}
\usepackage[unicode,hidelinks,bookmarks=false]{hyperref}
\newcommand{\ie}{i.e.\ }
\newcommand{\cf}{cf.\ }
\newcommand{\resp}{respectively\ }
\newcommand{\Subsection}{\subsection}
\providecommand{\nobreakdash}{\nobreak\hbox{-}}
\setlength{\emergencystretch}{2em}
\multlinegap0pt
% ---------------------------------------------------------------------------
% Editorial AMS wrapper prelude. The theorem environments and the mathematical macros below are
% transcribed from the paper's own preamble (cached-originals/source0.tex lines 32--99); the AMS amsart
% class, the named format packages of the pinned TeX Live runtime and this editorial formatting are not
% part of the author's text. No mathematics is changed.
% ---------------------------------------------------------------------------
\usepackage[nameinlink,noabbrev,nosort]{cleveref}
\newtheorem{thm}{Theorem}[section]
\newtheorem{lemma}[thm]{Lemma}
\newtheorem{ex}[thm]{Example}
\newtheorem{cl}{Claim}
\newtheorem{prop}[thm]{Proposition}
\newtheorem{con}[thm]{Conjecture}
\newtheorem{fact}{Fact}
\newtheorem{ques}[thm]{Question}
\newtheorem{cor}[thm]{Corollary}
\newtheorem{defn}[thm]{Definition}
\newtheorem{rem}[thm]{Remark}
\newtheorem{step}{Step}

\Crefname{thm}{Theorem}{Theorems}
\Crefname{lemma}{Lemma}{Lemmas}
\Crefname{prop}{Proposition}{Propositions}
\Crefname{subsection}{Subsection}{Subsections}

\newcommand{\norm}[1]{\left\Vert #1\right\Vert}
\newcommand{\nnorm}[1]{\lvert\!|\!| #1|\!|\!\rvert}

\renewcommand{\thefootnote}{$\dagger$}

\def \la {\longrightarrow}
\def \ul {\underline}
\def \ol {\overline}
\def \N {\mathbb N}
\def \C {\mathbb C}
\def \Z {\mathbb Z}
\def \R {\mathbb R}
\def \Q {\mathbb Q}
\def \E {\mathbb E}

\def\A {\mathcal A}
\def \M {\mathcal M}
\def\W {\mathcal W}
\def\U {\mathcal U}
\def\V {\mathcal V}
\def\bu {{\bf U}}
\def\I {\mathcal I}
\def\J {\mathcal J}
\def\B {\mathcal B}
\def\P {\mathcal P}

\def \X {\mathcal{X}}
\def \Y {\mathcal{Y}}
\def \ZZ {\mathcal{Z}}
\def \O {\mathcal{O}}
\def \t {\mathcal{T}}

\def \a {\alpha }
\def \b {\beta}
\def \ep {\epsilon}
\def \d {\delta}
\def \D {\Delta}

\def\lk {\left}
\def\re {\right}

\def\FS {{\bf FS}}
\def\FP {{\bf FP}}

\def\AP {{\bf AP}}
\def\GP {{\bf GP}}
\parskip 1.0ex
\numberwithin{equation}{section}
\def\htop{h_{\rm top}}


\begin{document}
\title{Pointwise convergence of double ergodic averages along certain non-polynomial sequences}
\author{Rongzhong Xiao}
\date{Source: arXiv:2510.17127v2; distributed under CC BY 4.0}
\begin{abstract}
	Fix $c\in (1,2)$. Let $\alpha$ and $\beta$ be two non-zero real numbers. It is shown that for any measure preserving system $(X,\X,\mu,T)$ and any $f,g\in L^{\infty}(\mu)$, the limit
	\begin{equation*}
	\lim_{N\to\infty}\frac{1}{N}\sum_{n=1}^{N}f(T^{\lfloor \alpha n^c \rfloor}x)g(T^{\lfloor \beta n^c \rfloor}x)
	\end{equation*}
	exists for $\mu$-a.e. $x\in X$.
\end{abstract}


\maketitle
\noindent\emph{Source, version and licence.} \emph{Pointwise convergence of double ergodic averages along certain non-polynomial sequences}, by Rongzhong Xiao. The source of this editable unit is the e-print of \textbf{version v2} of arXiv:2510.17127, https://arxiv.org/abs/2510.17127v2, whose material is distributed under the Creative Commons Attribution 4.0 International licence, http://creativecommons.org/licenses/by/4.0/, as recorded in the arXiv licence metadata for this paper and shown on that abstract page. The whole of the body below is version v2, the newest version of the paper. Full rights notice: licenses/arxiv-2510.17127-CC-BY-4.0.txt.\par\smallskip
\noindent\textit{Editorial scope.} Counted unit: original Proposition 3.2 of the article, the statement carried by the source environment \texttt{prop} with cleveref label \texttt{prop3:scalar}, cached-originals/source0.tex lines 1008--1018, printed as \textbf{Proposition 3.2} exactly as in the article. It is delivered together with the author’s own complete proof as the single primary proof body, source0.tex lines 1019--1146 (128 lines): the \texttt{proof} environment that immediately follows the statement and carries out the whole spectral reduction. No proof marker is added and no mathematics is written, repaired or re-derived; the author’s notation and the printed slips of the source are kept literal. Necessary complete context is retained uncounted: source0.tex lines 925--955, Theorem 3.1, the main theorem of the paper whose proof this proposition serves and whose hypotheses the counted statement refers to, together with its display of the two non-polynomial sequences. The article’s own numbering is reproduced by explicit counter settings: the counter \texttt{thm}, declared by \texttt{\textbackslash newtheorem\{thm\}\{Theorem\}[section]} and shared by lemma, example, proposition, conjecture, question, corollary, definition and remark, is set before each retained passage, so the counted statement prints as Proposition 3.2 with the number the article gives it, and the main theorem prints as Theorem 3.1. Every \texttt{\textbackslash ref}, \texttt{\textbackslash eqref} and \texttt{\textbackslash Cref} inside every retained passage resolves inside this delivered file, and the single citation that occurs in the retained passages, \texttt{Folland2016}, is delivered as the author’s own bibliography entry transcribed verbatim from the article’s own \texttt{thebibliography} with the author’s own printed number. Because the article marks its labels with cleveref’s optional type argument (\texttt{\textbackslash label[prop]\{prop3:scalar\}}), the wrapper loads the same named \texttt{cleveref} package with the article’s own options and transcribes the article’s own \texttt{\textbackslash Crefname} declarations, so that those labels and the \texttt{\textbackslash Cref} commands resolve and print exactly as in the source. The retained passages contain no figure environment and no graphics-inclusion command, so no artwork is redistributed. The source’s own class is AMS \texttt{amsart} (\texttt{\textbackslash documentclass[reqno,12pt]\{amsart\}}); the e-print ships no class or style file, so no publisher class is shipped, used or redistributed, and the wrapper is the same AMS \texttt{amsart} class with the paper’s own macros and theorem declarations transcribed from its preamble. Every declared body of this unit is an exact, unmodified span of source0.tex: no body is a bounded passage comparison, \texttt{omit} was not used anywhere, and consequently no fidelity receipt is asserted in delivery-evidence.json. The complete concatenated original source is bundled as sources/187389/source0.tex. Omitted from this editable unit and therefore absent from it are source0.tex lines 1--113; 122--924; 956--1007; 1147--3356; 3362--3403, that is the article’s own preamble and title block, the introduction with the quoted results of the literature, the transference machinery of Section 2, the remaining propositions and lemmas of Section 3 with their long proofs, and the remaining bibliography entries, all of which remain present in the bundled source but are not delivered here. The AMS \texttt{amsart} wrapper, the named format packages of the pinned TeX Live runtime and this editorial source, licence and scope paragraph are editorial formatting only.\par\smallskip
\setcounter{section}{3}
\setcounter{equation}{0}
\setcounter{cl}{0}
\setcounter{fact}{0}
\setcounter{step}{0}
\setcounter{thm}{0}
\begin{thm}\label{thm3-1}
Let $\rho,\theta\in(0,1)$, and fix
$\eta\in C_c^\infty((0,\infty);\C)$ and
$\omega\in C_c^\infty((1/2,2);\C)$. Set
\begin{equation}
 \sigma=\frac{1-\theta}{2},\qquad
 d_\theta=1+\left\lceil\frac2\sigma\right\rceil,
 \qquad P=2^{d_\theta}.
 \label{eq3-add-parameters}
\end{equation}
For an integer $N\ge2$ and a real number $M\in[1,N^\theta]$, define
$K_{N,M}:\R\to\C$ by
\begin{equation}
 K_{N,M}(u)=
 \begin{cases}
 \displaystyle\eta(u)\sum_{m\in\Z}\omega(m/M)e(-mNu^\rho),&u>0,\\[1mm]
 0,&u\le0.
 \end{cases}
 \label{eq3-add-kernel}
\end{equation}
For every measurable flow ${\bf Y}=(Y,\Y,\nu,(S^t)_{t\in\R})$,
every $\alpha,\beta\in\R$, every $R>0$, and all $1$-bounded
$F,G\in L^\infty(\nu)$,
\begin{equation}
 \left\|\int_\R K_{N,M}(u)
 F(S^{\alpha Ru})G(S^{\beta Ru})\,du\right\|_2^2
 \lesssim_{\rho,\theta,\eta,\omega}N^{-\sigma/P}.
 \label{eq3-add-main-bound}
\end{equation}
The implicit constant is independent of $N,M,R,\alpha,\beta,{\bf Y},F,G$.
\end{thm}

\setcounter{section}{3}
\setcounter{equation}{4}
\setcounter{cl}{0}
\setcounter{fact}{0}
\setcounter{step}{0}
\setcounter{thm}{1}
\begin{prop}\label[prop]{prop3:scalar}
Under the assumptions of \Cref{thm3-1}, for every real
$H\in(0,1]$,
\begin{equation}
 \left\|\int_\R K_{N,M}(u)
 F(S^{\alpha Ru})G(S^{\beta Ru})\,du\right\|_2^2
 \lesssim_{|I_0|}\frac1H\int_0^H
                \sup_{\xi\in\R}|\mathcal A_h(\xi)|\,dh.
 \label{eq3-add-spectral-reduction}
\end{equation}
\end{prop}

\setcounter{section}{3}
\setcounter{equation}{5}
\setcounter{cl}{0}
\setcounter{fact}{0}
\setcounter{step}{0}
\setcounter{thm}{2}
\begin{proof}
Choose measurable representatives of $F,G$ bounded by one everywhere,
and put
\[
 V(u)(y)=K_{N,M}(u)F(S^{\alpha Ru}y)G(S^{\beta Ru}y).
\]
Joint measurability of ${\bf Y}$, together with separability of $L^2(\nu)$, implies that $u\mapsto V(u)$ is strongly measurable. Since $$\|V(u)\|_2\le|K_{N,M}(u)|$$ and $K_{N,M}\in C_c^\infty(\R)$,\[
 V\in L^1(\R;L^2(\nu))\cap L^2(\R;L^2(\nu)).
\]
%Changing representatives does not change its integral, by measure
%preservation at each fixed time and Fubini's theorem.

Define
$$V_H(u)=H^{-1}\int_0^H V(u+s)\,ds$$ for almost every $u\in\R$. Then
$$\int_\R V_H(u)\,du=\frac1H\int_0^H\int_\R V(u+s)\,du\,ds=\int_\R V(u)\,du.$$
The support of $V_H$ is contained in $J_H=[a_0-H,b_0]$, whose length is
$$|J_H|=|I_0|+H\le |I_0|+1.$$ 
Note that 
\begin{align*}
& \left\|\int_\R V\right\|_2=\sup_{\substack{W\in L^2(\nu)\\ \|W\|_2=1}} \left|\left\langle\int_{J_H}V_H(u)\,du,W\right\rangle\right|
\\ & \hspace{3cm}\le\sup_{\substack{W\in L^2(\nu)\\ \|W\|_2=1}}
\int_{J_H}|\langle V_H(u),W\rangle|\,du
\le\int_{J_H}\|V_H(u)\|_2\,du.
\end{align*}
Therefore,
\begin{align*}
 &\left\|\int_\R V\right\|_2^2
 \le \left(\int_{J_H}\|V_H(u)\|_2\,du\right)^2
 \le (|I_0|+1)\int_\R\|V_H(u)\|_2^2\,du\\
 &\hspace{3cm}=\frac{|I_0|+1}{H^2}\int_0^H\int_0^H\int_\R
       \langle V(u+s),V(u+s')\rangle\,du\,ds\,ds'.
\end{align*}
%The first inequality is the triangle inequality for Bochner integrals. More explicitly, for every $W\in L^2(\nu)$ with $\|W\|_2=1$,
%\[ 
% \le\int_{J_H}|\langle V_H(u),W\rangle|\,du
% \le\int_{J_H}\|V_H(u)\|_2\,du,
%\]
%and taking the supremum over $W$ gives the first inequality. The second inequality is scalar Cauchy--Schwarz. The final equality follows by expanding the Bochner integral; Fubini is justified by
%$|\langle V(u+s),V(u+s')\rangle|\le |K_{N,M}(u+s)||K_{N,M}(u+s')|$.

Translating $u$ by $s'$ in the innermost integral gives
\[
 \int_\R\langle V(u+s),V(u+s')\rangle\,du=C(s-s'),
 \qquad
 C(h):=\int_\R\langle V(u+h),V(u)\rangle\,du.
\]
For every integrable scalar function $\Phi$ on $[-H,H]$,
\[
 \int_0^H\int_0^H\Phi(s-s')\,ds\,ds'
 =\int_{-H}^H(H-|h|)\Phi(h)\,dh.
\]
Indeed, set $h=s-s'$ and $t=s'$. The square $[0,H]^2$ becomes
\[
 \{(h,t):-H\le h\le H,\ 0\le t\le H,\ 0\le t+h\le H\}.
\]
For fixed $h$, the allowed $t$-interval is
$[\max(0,-h),\min(H,H-h)]$, whose length is $H-|h|$. Applying the above arguments with $\Phi=C$ and using $H-|h|\le H$ gives
\[
 \left\|\int_\R V\right\|_2^2
 \le\frac{|I_0|+1}{H}\int_{-H}^H|C(h)|\,dh.
\]

For a fixed $h\in\R$, let
$F_h=(F\circ S^{\alpha Rh})\overline F$ and
$G_h=(G\circ S^{\beta Rh})\overline G$.
Changing variables by $z=S^{\alpha Ru}y$ in the inner product gives
\begin{align*}
 \langle V(u+h),V(u)\rangle
 ={}&K_{N,M}(u+h)\overline{K_{N,M}(u)}
 \int_Y F_h(z)G_h(S^{(\beta-\alpha)Ru}z)\,d\nu(z).
\end{align*}
Set $U_uW=W\circ S^{(\beta-\alpha)Ru}$ on $L^2(\nu)$.
The last integral is $\langle U_uG_h,\overline{F_h}\rangle$.

We need to verify strong continuity of the action $U$ before applying the spectral theorem. Joint measurability of ${\bf Y}$ implies that
$t\mapsto\langle U_tW,Z\rangle$ is measurable for every $W,Z\in L^2(\nu)$. For
$W\in L^2(\nu)$ and $a\in C_c^\infty(\R)$, define 
\[
 W_a=\int_\R a(t)U_tW\,dt.
\]
The change of variables $r=t+s$ gives
\[
 U_sW_a=\int_\R a(r-s)U_rW\,dr,
\]
and hence
\[
 \|U_sW_a-W_a\|_2
 \le \|a(\cdot-s)-a\|_{L^{1}(\R)}\|W\|_2\longrightarrow0
 \qquad(s\to0).
\]
It remains to show that the span of the vectors $W_a$ is dense in $L^2(\nu)$. Suppose that
$Z$ is orthogonal to every such vector. Then
\[
 0=\langle W_a,Z\rangle
 =\int_\R a(t)\langle U_tW,Z\rangle\,dt
\]
for every $a\in C_c^\infty(\R)$. Thus
$t\mapsto\langle U_tW,Z\rangle$ vanishes almost everywhere for every $W$.
Choose a countable dense subset $\mathcal D\subset L^2(\nu)$ and a single
$t_0$ outside the union of the corresponding null sets. Then
$\langle U_{t_0}W,Z\rangle=0$ for every $W\in\mathcal D$, hence for every
$W\in L^2(\nu)$. Since $U_{t_0}$ is unitary, $Z=0$. Therefore, the span of the vectors $W_a$ is dense in $L^2(\nu)$. This proves the strong continuity of the action $U$.

By the spectral theorem (see \cite[Theorem~4.45]{Folland2016}), there is a
projection-valued measure $E$ on $\R$ such that
$$U_u=\int_\R e(u\xi)\,dE(\xi).$$ For $\varkappa\in L^1(\R)$, we have
\[
 \int_\R\varkappa(u)U_u\,du
 =\int_\R\left(\int_\R\varkappa(u)e(u\xi)\,du\right)dE(\xi).
\]
The operator on the left is defined strongly on each vector. The
spectral functional calculus bounds its norm by the supremum of the
scalar function in parentheses. Apply this to
$\varkappa(u)=K_{N,M}(u+h)\overline{K_{N,M}(u)}$ to obtain
\[
 \left|\int_\R\langle V(u+h),V(u)\rangle\,du\right|
 \le\sup_{\xi\in\R}|\mathcal A_h(\xi)|.
\]
Finally,
$\mathcal A_{-h}(\xi)=e(\xi h)\overline{\mathcal A_h(-\xi)}$,
so the two half-intervals in $h$ give equal suprema.
%Continuity in $\xi$ allows the supremum to be taken over $\Q$;
%it is therefore measurable in $h$. Its bound by
%$\|K_{N,M}\|_{L^2(\R)}^2$ also gives integrability.
This proves \eqref{eq3-add-spectral-reduction}.
%including the case
%$\alpha=\beta$, for which $U$ is the identity group.
\end{proof}

\setcounter{section}{0}
\setcounter{equation}{0}
\setcounter{cl}{0}
\setcounter{fact}{0}
\setcounter{step}{0}
\setcounter{thm}{0}
\begin{thebibliography}{99}
\bibitem[6]{Folland2016} G. B. Folland. \newblock \emph{A Course in Abstract Harmonic Analysis}, second edition. \newblock Textbooks in Mathematics, CRC Press, Boca Raton, FL, 2016. 
\end{thebibliography}
\end{document}
